%% chapter2.tex – Literature Survey
\chapter{Literature Survey}
\label{chap:litsurvey}

\section{Existing Coastal Monitoring Systems}

Global ocean monitoring has historically relied on moored buoys operated by national
agencies.
The U.S.\ National Data Buoy Center (NDBC) maintains hundreds of deep-water stations
that transmit wave height, period, and direction via satellite~\cite{ndbc2023}.
India's NIOT has deployed a network of Ocean Observation Buoys (OOBs) in the Arabian
Sea and Bay of Bengal, providing hourly wave spectra~\cite{niot2019buoys}.
While these systems are scientifically robust, they are designed for offshore deep water
(depths $>$100~m), cost several lakhs of rupees per unit, and require specialised ship
operations for deployment and retrieval.
None are positioned within 5~km of the Kerala near-shore zone that suffers from
\textit{Kallakkadal}.

Raj et al.~\cite{raj2018coastal} proposed a GPS drifter with accelerometry for shallow
coastal profiling in Goa, demonstrating that MEMS accelerometers can resolve wave
heights down to 0.2~m.
However, their prototype lacked onboard decision logic and GSM alerting.

\section{IoT in Disaster Management}

The convergence of low-cost microcontrollers, MEMS sensors, and cellular connectivity
has enabled a new generation of distributed environmental monitoring systems.
Postolache et al.~\cite{postolache2012iot} demonstrated a ZigBee-based flood sensor
network for Portuguese rivers, achieving 10-minute update cycles.
Suman and Jain~\cite{suman2015flood} used Arduino Uno with ultrasonic range finders and
a GSM shield to trigger SMS alerts when river levels exceeded safe thresholds; their
approach shares architectural similarities with WaveGuard~V4 but targets riverine rather
than marine contexts.
More recently, Mahapatra et al.~\cite{mahapatra2021esp32} replaced Arduino with ESP32,
exploiting the built-in WiFi to deliver real-time alerts to a web dashboard with less
than 500~ms latency.

\section{MEMS Accelerometry for Wave Sensing}

MEMS accelerometers have been validated for wave sensing in several studies.
Boore and Bommer~\cite{boore2005review} established the fundamental principles of
processing triaxial acceleration records; their spectral correction methods informed
our RMS-based threshold derivation.
Güven et al.~\cite{guven2020buoy} integrated an MPU6050 on a small buoy to study
harbour sloshing, reporting that a 50~Hz sample rate and $\pm$2g range are adequate for
wave periods of 6--25~s, matching the \textit{Kallakkadal} swell period of 10--20~s.
The 16\,384~LSB/g sensitivity corresponds to a resolution of $\approx$0.6~mg per LSB,
well below the 0.03~g \textsc{calm/moderate} threshold used in WaveGuard~V4.

\section{GSM and WiFi in Remote Alerting}

SIM800L-based alert systems have been widely studied for low-power IoT applications.
Al-Fuqaha et al.~\cite{alfuqaha2015iot} surveyed IoT protocols and confirmed that
GPRS-based SMS remains the most reliable alert channel in areas with poor broadband
coverage, which characterises the Kanayannur near-shore zone during monsoon conditions.
Taha et al.~\cite{taha2019gsm} measured SIM800L power transients during SMS
transmission, reporting 1.8--2.2~A peaks lasting approximately 500~ms---directly
informing WaveGuard~V4's decision to decouple the SIM800L power rail with a 1\,000~µF
bulk capacitor.

WiFi-based real-time dashboards for sensor data were studied by
Kodali et al.~\cite{kodali2016esp8266}, who demonstrated ESP8266 MQTT dashboards with
sub-200~ms latency. WaveGuard~V4 employs the successor ESP32 and REST/WebSocket
rather than MQTT, achieving comparable latency.

\section{Power Management for Marine IoT}

Prolonged field deployment requires careful energy budgeting.
Dondi et al.~\cite{dondi2008energy} modelled Li-ion discharge curves for remote
sensors, establishing that 18650 cells maintain $>$3.5~V under 150~mA load for
approximately 8~hours from a 2\,200~mAh starting charge---consistent with our
empirical 6--8~hour measurement.
The MT3608 boost converter chosen for WaveGuard~V4 is a common choice for 5~V MCU
rails from Li-ion sources, as evaluated by Tran et al.~\cite{tran2021boost}, who
reported $>$90\% efficiency at loads below 300~mA.

\section{Real-Time Backend Architectures}

FastAPI, introduced by Ramírez~\cite{fastapi2019}, provides asynchronous HTTP endpoints
with automatic OpenAPI documentation, making it suitable for sensor data ingestion at
high frequency.
SQLite is recommended for embedded backends where data rates do not exceed a few hundred
rows per second~\cite{owens2010sqlite}, which matches the WaveGuard~V4 ingestion rate
of 1~row/second.

\section{Related Coastal Early-Warning Projects}

The IOC-UNESCO Tsunami Warning System~\cite{ioc2006tsunami} provides the conceptual
template for multi-tier thresholding (green/amber/red), which WaveGuard~V4 adapts as
calm/moderate/high-wave.
Jayaratne et al.~\cite{jayaratne2016tsunami} and
Lauterjung et al.~\cite{lauterjung2010iotws} describe community-level siren networks
in Indonesia and Sri Lanka that complement satellite systems;
WaveGuard~V4 achieves the same last-mile delivery via SMS.

\section{Research Gap and WaveGuard Contribution}

Table~\ref{tab:litgap} summarises the gaps in existing work addressed by WaveGuard~V4.

\begin{table}[H]
\centering
\caption{Comparison of related systems with WaveGuard~V4}
\label{tab:litgap}
\begin{tabularx}{\textwidth}{X c c c c c}
\toprule
\textbf{System} & \textbf{Near-shore} & \textbf{MEMS} & \textbf{GSM} & \textbf{Dashboard} & \textbf{Cost} \\
\midrule
NDBC Buoy~\cite{ndbc2023}          & \ding{55} & \ding{55} & \ding{55} & \ding{51} & High   \\
NIOT OOB~\cite{niot2019buoys}      & \ding{55} & \ding{55} & \ding{55} & \ding{51} & High   \\
Raj et al.~\cite{raj2018coastal}   & \ding{51} & \ding{51} & \ding{55} & \ding{55} & Medium \\
Suman \& Jain~\cite{suman2015flood}& \ding{55} & \ding{55} & \ding{51} & \ding{55} & Low    \\
\textbf{WaveGuard~V4}              & \textbf{\ding{51}} & \textbf{\ding{51}} & \textbf{\ding{51}} & \textbf{\ding{51}} & \textbf{Low} \\
\bottomrule
\end{tabularx}
\end{table}

\noindent WaveGuard~V4 is the only system in the comparison that simultaneously achieves
near-shore deployment, MEMS-based wave sensing, GSM alerting, and a real-time web
dashboard at low cost, filling a clear gap in the existing literature.
